\documentclass[11pt]{article}

\usepackage[margin=0.85in]{geometry}
\usepackage{amsmath,amssymb}
\usepackage{enumitem}
\usepackage{xcolor}
\usepackage{tcolorbox}
\usepackage{fancyhdr}
\usepackage{lastpage}
\usepackage{hyperref}
\usepackage{microtype}

\definecolor{uwpurple}{HTML}{4B2E83}
\definecolor{uwgold}{HTML}{B7A57A}
\definecolor{softpurple}{HTML}{F4F1FA}

\hypersetup{
    colorlinks=true,
    linkcolor=uwpurple,
    urlcolor=uwpurple
}

\pagestyle{fancy}
\fancyhf{}
\lhead{\textbf{MATH 394: Probability I}}
\rhead{\textbf{Homework 5}}
\cfoot{\thepage\ of \pageref{LastPage}}

\setlength{\parindent}{0pt}
\setlength{\parskip}{0.45em}

\newcommand{\course}{MATH 394: Probability I}
\newcommand{\instructor}{Arman Jahangiri}
\newcommand{\department}{Department of Mathematics}
\newcommand{\university}{University of Washington}

\newtcolorbox{problemheader}{
    colback=softpurple,
    colframe=uwpurple,
    boxrule=0.8pt,
    arc=2mm,
    left=2mm,
    right=2mm,
    top=1mm,
    bottom=1mm
}

\newcommand{\source}[1]{%
    \vspace{0.15cm}
    {\small\textit{Source: Blitzstein and Hwang, \textit{Introduction to Probability}, Chapter 4, Problem #1.}}
}

\newcommand{\answerbox}{%
\begin{flushright}
\begin{tcolorbox}[
    width=0.36\textwidth,
    height=2cm,
    colback=white,
    colframe=uwpurple,
    boxrule=0.7pt,
    arc=1mm
]
{\small\textbf{Final Answer}}
\end{tcolorbox}
\end{flushright}
}

\newcommand{\partanswerbox}[1]{%
\begin{flushright}
\begin{tcolorbox}[
    width=0.36\textwidth,
    height=2cm,
    colback=white,
    colframe=uwpurple,
    boxrule=0.7pt,
    arc=1mm
]
{\small\textbf{Final Answer for Part #1}}
\end{tcolorbox}
\end{flushright}
}

\begin{document}

\begin{titlepage}
\centering
\vspace*{1.2cm}

{\Large \textbf{\university}}\\
{\large \department}

\vspace{1.5cm}

{\Huge \color{uwpurple}\textbf{Homework 5}}\\[0.25cm]
{\Large \textbf{\course}}\\[0.25cm]
{\large Instructor: Arman Jahangiri}

\vspace{1cm}

\begin{tcolorbox}[
    width=0.78\textwidth,
    colback=softpurple,
    colframe=uwpurple,
    boxrule=0.9pt,
    arc=3mm
]
\centering
\textbf{Submission:} Single PDF on Gradescope

\textbf{Deadline:} 10:00 PM Pacific Time on the listed due date
\end{tcolorbox}

\vspace{0.8cm}

\begin{tcolorbox}[
    width=0.86\textwidth,
    colback=white,
    colframe=uwgold,
    boxrule=0.8pt,
    arc=3mm
]
This homework is worth \textbf{50 points}.

Across the quarter, there are \textbf{eight homework assignments}, worth \textbf{400 points total}. Homework assignments together account for \textbf{40\%} of the final course grade. Further course policies can be seen in the \href{https://armanjg.github.io/MATH-394-Probability-1/}{\textbf{MATH 394 syllabus}.}
\end{tcolorbox}

\vfill
\end{titlepage}

\section*{Homework Policy}

\subsection*{Submission and Deadline}

Please:

\begin{itemize}[leftmargin=*]
    \item Submit homework as a single PDF on Gradescope by \textbf{10:00 PM (Pacific Time)} on the listed due date.
    \item Begin each problem on a new page, and ensure that pages are correctly tagged in Gradescope.
    \item Write solutions clearly and legibly.
    \item Typesetting solutions in LaTeX is encouraged, but not required.
\end{itemize}

\subsection*{Late Days}

Each student is allotted \textbf{six late days} for the quarter. A late day extends a homework deadline by up to 24 hours without penalty. For example, submitting an assignment anytime between 10:01 PM on the due date and 10:00 PM the following day counts as one late day.

The following rules apply:
\begin{itemize}[leftmargin=*]
   \item No explanation is required to use late days.
   \item At most \textbf{two late days} may be used on any single homework assignment.
   \item Late days may not be used on the final homework assignment or on any assignment for which solutions have already been released.
   \item It is the student's responsibility to keep track of how many late days have been used during the quarter.
\end{itemize}

Once all late days have been exhausted, additional late submissions will incur a penalty of \textbf{10\% per day}, up to a maximum deduction of \textbf{50\%}. Assignments submitted more than five days late, or after solutions have been released, will not be accepted.

\subsection*{Submission Issues and Technical Difficulties}

If a serious technical issue prevents a timely Gradescope submission, students may temporarily submit their assignment by email to the instructor at \href{mailto:armanjg@uw.edu}{armanjg@uw.edu}. In such cases:

\begin{itemize}[leftmargin=*]
    \item The submission must consist of a \textbf{single PDF file}. The email subject line must follow the format:
    \[
    \texttt{HW\# Submission - FULL NAME - STUDENT ID}.
    \]
    \item Once the Gradescope issue has been resolved, the student must upload the \emph{exact same file} to Gradescope as soon as possible.
    \item The student should clearly indicate in the Gradescope submission that the assignment was previously submitted by email before the deadline.
\end{itemize}

Email submissions are intended only for genuine technical emergencies and should not be used as a substitute for timely Gradescope submission.

\newpage

% ------------------------------------------------------------

\begin{problemheader}
\textbf{Problem 1. (10 points)}
\end{problemheader}
\source{18}

A fair coin is tossed repeatedly, until it has landed Heads at least once and has landed Tails at least once.

Find the expected number of tosses.

 
 
\textit{Hint}: Define \(T\) to be the total number of tosses required. After the first toss,
only one outcome remains to be observed. Write
\[
T=1+X,
\]
where \(X\) is the number of additional tosses needed to obtain the outcome
opposite to the first toss. $X$ follows a well-known distribution! So if you find that, the rest is to use the linearity of expectation.
 
 
 
\newpage
 

% ------------------------------------------------------------

\begin{problemheader}
\textbf{Problem 2. (10 points) }
\end{problemheader} 
\source{30}

\begin{enumerate}[label=(\alph*), leftmargin=*]
    \item Use LOTUS to show that for $X\sim \operatorname{Pois}(\lambda)$ and any function $g$,
    \[
    E\bigl(Xg(X)\bigr)=\lambda E\bigl(g(X+1)\bigr),
    \]
    assuming that both sides exist. This is called the \emph{Stein--Chen identity} for the Poisson.

    \item Find the third moment $E(X^3)$ for $X\sim\operatorname{Pois}(\lambda)$ by using the identity from part (a) and a bit of algebra to reduce the calculation to the fact that $X$ has mean $\lambda$ and variance $\lambda$.
\end{enumerate}
 
\newpage

% ------------------------------------------------------------


% ------------------------------------------------------------

% ------------------------------------------------------------

% ------------------------------------------------------------
% ------------------------------------------------------------

\begin{problemheader}
\textbf{Problem 3.  (30 points)}
\end{problemheader}
\source{54}

\textbf{Note.}
This problem introduces several fundamental ideas in data analysis, including
sampling, estimation, uncertainty, bias, and statistical precision. Pay
particular attention to this problem, specially if you'd like to smile, remembering this course, years later in a job that contains data collection/analysis. :)

A population consists of \(N\) people, with ID numbers \(1,\dots,N\). For each
person \(j\), let \(y_j\) be the value of some numerical variable, and define the
population mean by
\[
\bar y=\frac{1}{N}\sum_{j=1}^N y_j.
\]

A quantity such as \(\bar y\) is called a \emph{population parameter}. Once the
population is fixed, its value is fixed, although it may be unknown to us.
Population parameters are often denoted by symbols such as
\(\theta\), \(\mu\), or \(\bar y\).

In contrast, a quantity computed from a random sample is called a
\emph{sample statistic}. For example, if \(W_1,\dots,W_n\) are the values
observed in a random sample, then
\[
\bar W
=
\frac{1}{n}\sum_{r=1}^n W_r
\]
is a sample statistic. Since the selected sample is random, \(\bar W\) is a
random variable whose value may vary from one sample to another. More generally,
a statistic may be written as
\[
T=T(W_1,\dots,W_n),
\]
emphasizing that it is a function of the observed sample.

One of the primary goals of statistics is to use sample statistics to estimate
unknown population parameters.

Suppose a researcher wants to estimate the population parameter
\[
\theta:=\bar y,
\]
but it is not feasible to measure every person in the population (which is required for computing $\bar y$ based on its formula). The researcher
therefore selects a simple random sample of size \(n\), where \(n\ll N\),
without replacement. Thus, every subset of \(n\) people is equally likely to be
selected.

Let \(W_1,\dots,W_n\) denote the values observed for the selected people, where
the sampled people are listed in a uniformly random order, and let
\[
\bar W=\frac{1}{n}\sum_{r=1}^n W_r
\]
be the sample mean.

For each population member \(1\leq j\leq N\), define
\[
I_j=
\begin{cases}
1, & \text{if person \(j\) is included in the sample},\\
0, & \text{otherwise}.
\end{cases}
\]

Define the finite-population variance by
\[
S_y^2
=
\frac{1}{N-1}\sum_{j=1}^N (y_j-\bar y)^2.
\]

Answer the following questions.

\begin{enumerate}[label=(\alph*), leftmargin=*]

\item Explain carefully why \(y_1,\dots,y_N\) are not random variables, while
\(W_1,\dots,W_n\), \(I_1,\dots,I_N\), and \(\bar W\) are random variables.

\item Show by symmetry that, for every sample position \(r\),
\[
P(W_r=y_j)=\frac{1}{N},
\qquad j=1,\dots,N.
\]
Use this fact to prove that
\[
E(\bar W)=\bar y.
\]

\item Show that
\[
\bar W=\frac{1}{n}\sum_{j=1}^N I_jy_j.
\]
Then compute \(E(I_j)\) and use linearity of expectation to give a second proof
that
\[
E(\bar W)=\bar y.
\]

\item Since the sample contains exactly \(n\) people, show that
\[
\sum_{j=1}^N I_j=n
\]
with probability \(1\). Explain why this identity does not imply that the
individual indicators \(I_1,\dots,I_N\) are constants.



\item Compute \(\operatorname{Var}(I_j)\). Then use the indicator
representation of \(\bar W\) and the definition of variance to show that
\[
\operatorname{Var}(\bar W)
=
\frac{N-n}{Nn}S_y^2
=
\left(1-\frac{n}{N}\right)\frac{S_y^2}{n}.
\]
You may use the fact that, for \(i\neq j\),
\[
P(I_i=1 \text{ and } I_j=1)
=
\frac{n(n-1)}{N(N-1)}.
\]
The factor
\[
1-\frac{n}{N}
\]
is called the \emph{finite-population correction}.





 

\item Explain what the variance formula predicts in each of the following
cases:
\begin{enumerate}[label=(\roman*), leftmargin=2em]
    \item \(n=1\);
    \item \(n=N\);
    \item all population values are equal;
    \item \(n\) increases while the population remains fixed.
\end{enumerate}

\item Suppose instead that the researcher samples \(n\) people independently
\emph{with replacement}. Let \(\bar W_{\mathrm{rep}}\) be the resulting sample
mean. Show that
\[
\operatorname{Var}(\bar W_{\mathrm{rep}})
=
\frac{N-1}{N}\frac{S_y^2}{n}.
\]
Compare this with the variance under sampling without replacement. Why is
sampling without replacement more precise?

\item The bias and mean squared error of an estimator \(T\) for a population
parameter \(\theta\) are defined by
\[
\operatorname{Bias}(T)=E(T)-\theta
\]
and
\[
\operatorname{MSE}(T)=E\bigl[(T-\theta)^2\bigr].
\]
Show that
\[
\operatorname{MSE}(T)
=
\operatorname{Var}(T)+\operatorname{Bias}(T)^2.
\]
Use this identity to find \(\operatorname{Bias}(\bar W)\) and
\(\operatorname{MSE}(\bar W)\).

\item A researcher says:
\begin{quote}
``Because \(\bar W\) is unbiased, the sample mean must equal the population mean
for every possible sample.''
\end{quote}
Explain why this statement is false. What does unbiasedness actually mean?

\end{enumerate}
 

\end{document}